\chapter{Conclusions}

\framebox{Note:} this is very much a first draft.

This book has presented a number of tools and techniques for constructing
concurrent programs.  

We have advocated constructing concurrent programs around concurrent objects.
Each concurrent object can be developed and tested in isolation: thinking
about a single object at a time is much more straightforward than global
reasoning.  The client code can then use these concurrent objects much the
same as sequential code.  We have seen patterns, such as the bag-of-tasks
pattern, that directly use this technique.

We have seen various techniques for implementing a concurrent object: message
passing (using the client-server pattern); explicit locks and signals (either
SCL or JVM monitors); or semaphores.  Message passing can be more
straightforward, particularly for synchonisation objects, where a single
thread---the server---can decide how to proceed.  However, if out-and-out
performance is important, then JVM monitors tend to be faster.  If a
particular concurrent object turns out to be a bottleneck, then fine-grained
locking or sharding can prove effective. 

A concurrent object can be tested in isolation.  Typically, this involves
logging interactions with the object, and then analysing the log for
correctness.  In particular, for concurrent datatypes, this can be done within
the linearisability framework. 

It is often necessary to provide extra coordination between threads.  For
programs that proceed in rounds, this can be achieved using a barrier
synchronisation at the end of each round, possibly with extra synchronisations
between different stages of a round.  Alternatively, coordination can be
achieved via message passing, or a special purpose synchronisation object.

In the remainder of this chapter, we briefly discuss some other patterns for
concurrent programs. 

%%%%%%%%%%%%%%%%%%%%%%%%%%%%%%%%%%%%%%%%%%%%%%%%%%%%%%%

\section{Atomics}

In a concurrent program, particularly one with many threads, certain
components can act as a bottleneck.  For example, if a concurrent object uses
locking, then when one thread holds the lock, others have to wait.

An alternative approach is not to use locking.  However, how can we then avoid
races of the form we discussed back in Chapter~\ref{chap:intro}?  The answer
is that we need to make use of additional operations provided by the hardware.
In particular, we can use an operation |compareAndSet| (CAS) on a memory
location, which can be described in code as follows, where |current| is the
location in question, and ``|atomically|'' is intended to indicate that the
operation is performed atomically.
%
\begin{scala}
  def compareAndSet(expected: Int, newValue: Int): Boolean = atomically{
    if(current == expected){ current = newValue; true } else false
  }
\end{scala}
%
The operation takes parameters |expected|, which represents the value the
client code expects the location to have, and |newValue|, which is a desired
new value.  If the current value is indeed the expected value, then it is
updated to |newValue|, and the operation returns |true|; otherwise the
operation returns |false|.

In Scala (and Java) these operations are encapsulated in objects from the
package |java.util.concurrent.atomic|.  In particular, each object from the
class |AtomicInteger| encapsulates an |Int|, and each object from the class
|AtomicReference[A]| encapsulates a reference to an object of type~|A|.  These
objects provide the |compareAndSet| operation, a |get| operation to get the
current value, and more.

For example, we could increment an |AtomicInteger| using the following
function. 
%
\begin{scala} 
  def increment(x: AtomicInteger) = {
    var done = false
    while(!done){
      val oldX = x.get
      done = x.compareAndSet(oldX, oldX+1)
    }
  }
\end{scala}
%
This reads the current value of~|x| into |oldX|, and then tries to use a CAS
to change the value of~|x| from |oldX| to~|oldX+1|.  If this succeeds, it is
done.  Otherwise, some other thread has changed the value of~|x| in the
meantime, so that it no longer equals |oldX|; the operation retries in this
case.

An operation is said to be \emph{lock-free}\index{lock-free} if, whenever
there is a pending operation call, some such call eventually returns.  The
above |increment| operation is lock-free: if the CAS of one thread fails, some
other thread must have updated~|x|, and that other thread can return. 

%%  On the other hand, an operation that uses a lock is (unsurprisingly) not
%% lock-free: if one thread obtains the lock and is permanently suspended,
%% then other threads will be blocked attempting to obtain the lock, and so no
%% operation call will return.  (In practice, threads don't get permanently
%% suspended; however, they can be suspended for a period, and this can block
%% other threads.)


As a larger example, Figure~\ref{fig:lock-free-stack} give the definition of a
lock-free stack.  The implementation is built around a linked list, built from
|Node|s.  It maintains a reference |top| to the first node in that list.  The
|push| operation creates a new |Node|, and sets it to point to the current
|top|.  It then attempts to update |top| using a CAS to point to the new node.
If this fails, some other thread has updated |top| in the meantime, so it no
longer equals |oldTop|, so this thread retries.  If the |pop| operation finds
that |top| holds |null|, the stack is empty, so it returns |None|.  Otherwise,
it tries to update |top| using a CAS to point to the next node; if this fails,
again some other thread has updated |top|, so this thread retries.


\begin{figure}
\begin{scala}
class LockFreeStack[A] extends TotalStack[A]{
  /** Nodes in the linked list. */
  private class Node(val value: A){
    var next: Node = null
  }

  /** Atomic reference to the top of the stack. */
  private val top = new AtomicReference[Node](null)

  /** Push £value£ onto the stack. */
  def push(value: A) = {
    val node = new Node(value); var done = false
    while(!done){
      val oldTop = top.get; node.next = oldTop
      // Try to add node to the stack.
      done = top.compareAndSet(oldTop, node) 
    } 
  }

  /** Pop a value from the stack. Return £None£ if the stack is empty.  */ 
  def pop: Option[A] = {
    var result: Option[A] = None; var done = false
    while(!done){
      val oldTop = top.get
      if(oldTop == null) done = true // Empty stack; return £None£.
      else{
        val newTop = oldTop.next
        // Try to remove £oldTop£ from the list.
        if(top.compareAndSet(oldTop, newTop)){
          result = Some(oldTop.value); done = true
        }
      }
    }
    result
  }
}
\end{scala}
\caption{A lock-free total stack (implementing the {\scalastyle TotalStack}
  trait from Figure~\protect\ref{fig:immutable-stack}).}
\label{fig:lock-free-stack}
\end{figure}

%%%%%

\begin{instruction}
Make sure you understand the details of the implementation, including why it
is necessary to retry if a CAS fails. 
\end{instruction}

An operation that uses a lock is (unsurprisingly) not lock-free: if one thread
obtains the lock and is permanently suspended, then other threads will be
blocked attempting to obtain the lock, and so no operation call will return.
In practice, threads don't get permanently suspended; however, they can be
suspended for a period, and this can block other threads.  As we saw in
Section~\ref{sec:Amdahl}, when we discussed Amdahl's Law, this can mean that
the performance improves only slightly as the number of processors increases.
On the other hand, if an operation is lock-free, it is likely to scale better:
while a particular thread might have to retry repeatedly because of actions by
other threads, the delay is likely to be less than that caused by having to
wait to obtain a lock.  Having said that, correctly implementing lock-free
objects is considerably harder than using locks. 

%  spin lock??

%%%%%%%%%%%%%%%%%%%%%%%%%%%%%%%%%%%%%%%%%%%%%%%%%%%%%%%

\section{Futures}

A \emph{future} is a value that is computed in parallel with the main
computation, to be used at some time in the future.  For example:
%
\begin{scala}
  val x = new Future(...) // some lengthy computation. 
  ...
  val y = f(x())
\end{scala}
%
The computation of the value, encapsulated in~|x|, is started in the first
line.  The expression \SCALA{x()} (calling the |apply| method on~|x|; see
Section~\ref{sec:scala-misc}) returns the result of the computation, waiting
until it is complete if necessary.  (Of course, it is still necessary to avoid
race conditions.)

Figure~\ref{fig:future} gives an implementation of futures using a JVM
monitor.  It forks a separate thread to calculate the value required; a call
of the |apply| method has to wait until that calculation is complete.
%
\begin{instruction}
Study the details of the implementation. 
\end{instruction}

%%%%%

\begin{figure}
\begin{scala}
class Future[A](exp: => A){
  /** Variable to hold the result. */
  private var result = null.asInstanceOf[A]

  /** Is the result valid? */
  private var done = false

  /** Server thread to calculate the result. */
  private def server = thread{
    synchronized{ result = exp; done = true; notifyAll() }
  }

  fork(server)

  /** Get the value of £exp£. */
  def apply(): A = synchronized{ while(!done) wait(); result }
}
\end{scala}
\caption{An implementation of futures using a JVM monitor.}
\label{fig:future}
\end{figure}


%%%%%%%%%%%%%%%%%%%%%%%%%%%%%%%%%%%%%%%%%%%%%%%%%%%%%%%

\section{Competition parallel}

The idea of competition parallel\index{competition parallel} is to choose two
(or more) different algorithms for a problem, and to run them independently in
parallel; once one finishes, the others can be interrupted.

This works well on the boolean satisfiability problem.  Here, one is given a
boolean formula, such as $(b_1 \lor b_3) \land (\lnot b_1 \lor \lnot b_3)
\land (\lnot b_1 \lor b_2 \lor b_3)$, and one has to find if there are values
for the boolean variables that makes the formula true (e.g.,~$b_1 = b_2 =
true$,\, $b_3 = false$ in the above example).  The problem is NP-complete.
However, there are a number of SAT-solving algorithms that work well in many
cases.  But all known algorithms have cases where they perform badly, and
different algorithms perform badly on different cases.  Therefore running two
algorithms in competition with one another can give better average results
than parallelising a single algorithm.

%% Competition parallel can also be used effectively when searching, for example
%% in a breadth-first or depth-first search.  Such searches can be done in two
%% ways.  In a \emph{graph search}, the algorithm keeps track of previously
%% visited nodes, to avoid repeating work; this is the role of the |seen| set in
%% Figure~\ref{fig:bfs-seq}.  In a \emph{tree search}, the algorithm does not keep
%% track of those nodes.


I have used competition parallel effectively within linearisability testing.
Recall that here one has to test whether a given concurrent history is
linearisable with respect to a sequential specification object.  This is done
by a depth-first search.  Such searches can be done in two ways.  In a
\emph{graph search}, the algorithm keeps track of previously visited nodes, to
avoid repeating work (this is the role of the |seen| set in the breadth-first
search of Figure~\ref{fig:bfs-seq}).  In a \emph{tree search}, the algorithm
does not keep track of those nodes.  It turns out that, for linearisabilty
testing, the tree search works better for most histories, because it avoids
the additional cost of maintaining the set of nodes, and normally doesn't
repeat too much work.  However, the tree search sometimes hits a very bad
case, where the search does not terminate within a reasonable time, because it
repeatedly expands the same nodes.  Running these two algorithms in
competition parallel works well: for most histories, the tree search finishes
quickly; but when the tree search hits a bad case, the graph search still
finishes quite quickly.

%%%%%%%%%%%%%%%%%%%%%%%%%%%%%%%%%%%%%%%%%%%%%%%%%%%%%%%
